% oss-plantilla: oss-2026.11
% Preámbulo corporativo fijo. No editar en sd-NN.tex: exportar_latex.py verificar lo compara byte a byte.
% Para cambiar el formato, editar esta plantilla u oss.sty y subir la versión de la primera línea.
\PassOptionsToPackage{unicode}{hyperref}
\PassOptionsToPackage{hyphens}{url}
\documentclass[11pt]{article}
\usepackage{amsmath,amssymb}
\usepackage{iftex}
\ifPDFTeX
  % Prism y otros editores en la nube: pdfLaTeX con Helvetica como familia por defecto.
  \usepackage[T1]{fontenc}
  \usepackage[utf8]{inputenc}
  \usepackage{textcomp}
  \usepackage[scaled=0.92]{helvet}
  \renewcommand{\familydefault}{\sfdefault}
  \DeclareUnicodeCharacter{2264}{\ensuremath{\leq}}
  \DeclareUnicodeCharacter{2265}{\ensuremath{\geq}}
  \DeclareUnicodeCharacter{2260}{\ensuremath{\neq}}
  \DeclareUnicodeCharacter{2248}{\ensuremath{\approx}}
  \DeclareUnicodeCharacter{2190}{\ensuremath{\leftarrow}}
  \DeclareUnicodeCharacter{2192}{\ensuremath{\rightarrow}}
  \DeclareUnicodeCharacter{25CF}{\ensuremath{\bullet}}
  \DeclareUnicodeCharacter{2713}{\ensuremath{\checkmark}}
\else
  % XeLaTeX/LuaLaTeX (compilación oficial local): DejaVu Sans.
  \usepackage{fontspec}
  \defaultfontfeatures{Ligatures=TeX,Scale=MatchLowercase}
  \setmainfont{DejaVu Sans}
  \setsansfont{DejaVu Sans}
  \setmonofont{DejaVu Sans Mono}
\fi
\usepackage[letterpaper,margin=20mm]{geometry}
\usepackage{microtype}
\usepackage{parskip}
\usepackage{xcolor}
\usepackage{longtable,booktabs,array}
\usepackage{calc}
\usepackage{etoolbox}
\makeatletter
\patchcmd\longtable{\par}{\if@noskipsec\mbox{}\fi\par}{}{}
\makeatother
\usepackage{footnotehyper}
\makesavenoteenv{longtable}
\usepackage{graphicx}
\makeatletter
\newsavebox\pandoc@box
\newcommand*\pandocbounded[1]{%
  \sbox\pandoc@box{#1}%
  \Gscale@div\@tempa{\textheight}{\dimexpr\ht\pandoc@box+\dp\pandoc@box\relax}%
  \Gscale@div\@tempb{\linewidth}{\wd\pandoc@box}%
  \ifdim\@tempb\p@<\@tempa\p@\let\@tempa\@tempb\fi
  \ifdim\@tempa\p@<\p@\scalebox{\@tempa}{\usebox\pandoc@box}%
  \else\usebox{\pandoc@box}%
  \fi}
\def\maxwidth{\ifdim\Gin@nat@width>\linewidth\linewidth\else\Gin@nat@width\fi}
\def\maxheight{\ifdim\Gin@nat@height>\textheight\textheight\else\Gin@nat@height\fi}
\makeatother
\setkeys{Gin}{width=\maxwidth,height=\maxheight,keepaspectratio}
\makeatletter
\def\fps@figure{htbp}
\makeatother
\usepackage[normalem]{ulem}
\providecommand{\st}[1]{\sout{#1}}
\usepackage{fancyvrb}
\setlength{\emergencystretch}{3em}
\providecommand{\tightlist}{%
  \setlength{\itemsep}{0pt}\setlength{\parskip}{0pt}}
\providecommand{\passthrough}[1]{#1}
\setcounter{secnumdepth}{-\maxdimen}
\ifPDFTeX
  \usepackage[spanish]{babel}
\else
  \usepackage[bidi=default]{babel}
  \babelprovide[main,import]{spanish}
\fi
\let\LanguageShortHands\languageshorthands
\def\languageshorthands#1{}
\usepackage{bookmark}
\usepackage{xurl}
\urlstyle{same}
\hypersetup{pdflang={es-CL},pdfcreator={Only Simple Solutions - exportar_latex.py}}
\usepackage{oss}
\begin{document}

\ossCover{Plan de calidad}{Subdocumento 9}
\tableofcontents
\clearpage

\renewcommand{\thefigure}{9.\arabic{figure}}

\section{Capítulo 9 · Introducción al Plan de calidad}

Este capítulo establece cómo Only Simple Solutions planificará, asegurará y medirá la calidad de los productos y del trabajo que los genera durante las dos etapas de implementación y la operación contractual. Su propósito es que cada incremento y entregable de Ancoa cumpla los requisitos aprobados y llegue a sus hitos con evidencia revisable. El plan conecta las prácticas de desarrollo y DevSecOps del Subdocumento 6 con la arquitectura, los riesgos y el plan de trabajo de los Subdocumentos 4, 8 y 7. El Formulario T-13 desarrollará el plan y los casos de prueba; el T-17 detallará el protocolo de aceptación.

La calidad se observará tanto en el producto como en el proceso. En Ancoa, esto incluye verificar las reglas de venta, inventario y crédito, además de preservar la separación entre Retail y el Emisor fiscalizado en datos, accesos e integraciones (Multitiendas Ancoa S.A., 2026, Caso 09, cap. 9, p. 25). Las mediciones y revisiones producirán registros por versión, ambiente y requisito. Las campañas y sus registros se ejecutarán durante el proyecto; este plan define su diseño y condiciones de aceptación.

\section{9.1 Plan de Calidad}

El plan convierte los compromisos de las Bases en criterios de trabajo y mediciones para requisitos, arquitectura, código, datos, pruebas y entregas. La calidad es el grado en que el producto satisface los requisitos acordados; el grado describe el conjunto de prestaciones y capacidades elegido para el producto. Por ello, una decisión de alcance puede definir qué funciones se incluyen, pero no reduce la obligación de que cada función comprometida opere conforme a sus criterios. El plan cubre la calidad del proceso, que se controla con prácticas repetibles y revisiones, y la calidad del producto, que se verifica mediante pruebas, inspecciones y métricas. Prevención y evaluación se incorporan al trabajo; la corrección atiende defectos detectados antes o después de una entrega.

Para esta gestión, la conformidad se revisa frente a los requisitos aprobados, los estándares de trabajo documentados y las expectativas profesionales que puedan afectar seguridad, integridad o continuidad. Una expectativa implícita no crea por sí sola una obligación nueva. Se registra como supuesto o riesgo, se asigna a un responsable y se resuelve mediante el mecanismo de decisión o control de cambios antes de incorporarla a una línea base. El grado se expresa en las prestaciones y funciones escogidas; una vez comprometida una prestación, su calidad se juzga por el criterio aprobado para ella. Las revisiones del proceso comprueban que se aplicaron prácticas definidas, mientras las pruebas e inspecciones aportan evidencia sobre el producto; aprobar una revisión de proceso no demuestra por sí solo la conformidad del producto.

En Retail, elegir qué canales y perfiles se entregan en cada etapa es una decisión de grado delimitada por el alcance; cobrar el precio vigente y devolver un estado de venta consistente es calidad de las prestaciones comprometidas. En el Emisor, acelerar la evaluación crediticia no autoriza omitir autorización, consentimiento o evidencia precontractual. El resultado esperado debe reunir desempeño y cumplimiento en la misma operación.

La revisión inicial de expectativas examinará la conservación de una venta confirmada ante interrupción, la vigencia y fuente de la existencia mostrada, la ausencia de cobros duplicados ante reintentos, el estado seguro de una evaluación crediticia incompleta y el rastro de acceso a cartera. Arquitectura, Seguridad y el dueño del dominio contrastarán cada expectativa con las Bases y el Caso. Si ya existe una obligación, vincularán su requisito y prueba; si aparece una condición no expresada, la registrarán como riesgo o decisión de alcance para la Contraparte. El registro distinguirá fundamento contractual, supuesto y decisión, sin introducir obligaciones tácitas como hechos acordados.

\subsection{9.1.1 Marco de calidad y características del producto}

Evaluaremos cada entrega de Ancoa en las ocho categorías contractuales: funcionalidad, desempeño, compatibilidad, usabilidad, fiabilidad, seguridad, mantenibilidad y portabilidad. Para convertirlas en criterios medibles utilizaremos ISO/IEC 25010:2023, cuyo modelo actual de calidad del producto organiza nueve características. Su vocabulario ayuda a especificar pruebas y métricas sin sustituir los títulos de las Bases: capacidad de interacción desarrolla la antigua dimensión de usabilidad, flexibilidad desarrolla portabilidad y seguridad de funcionamiento añade una perspectiva para las fallas con consecuencias operativas. El Artículo 24 exige criterios de calidad y controles DevSecOps; el Formulario T-7 pide además tratar madurez y umbrales bloqueantes de complejidad y acoplamiento. Este plan responde a cada obligación en las secciones siguientes.

En Ancoa, la adecuación funcional comprobará los flujos de venta, inventario y crédito; la eficiencia de desempeño se medirá bajo los perfiles y cargas establecidos en las Bases; la compatibilidad cubrirá interfaces con las plataformas que permanecen; la capacidad de interacción atenderá los puestos de tienda y los procesos de atención; la fiabilidad cubrirá disponibilidad y recuperación; la seguridad verificará acceso, integridad y segregación entre Retail y Emisor; la mantenibilidad considerará pruebas y estructura del código; la flexibilidad observará adaptación y escalamiento; y la seguridad de funcionamiento evaluará el comportamiento seguro ante fallas que puedan afectar personas, bienes u operación. Los objetivos numéricos se derivan de las Bases o de decisiones documentadas del proyecto, no del modelo ISO por sí solo.

Seleccionamos el nivel 3 (Definido) de CMMI v3.0 como meta de madurez para los procesos que aplicaremos en Ancoa. Esta meta es adecuada para dos etapas solapadas y una operación prolongada: requisitos, planificación, verificación, configuración y mejora deben seguir procedimientos comunes y adaptaciones controladas, con registros que permitan repetirlos entre equipos. En el primer mes, el Líder de Calidad y el Jefe de Proyecto levantarán una línea base de esas prácticas; en cada cierre de intervalo de SD06 comprobarán procedimientos, revisiones y registros, y asignarán dueño y fecha a las brechas. Antes de certificar cada etapa revisarán si las prácticas previstas están implantadas y qué acciones siguen pendientes. El nivel 3 es una meta elegida por Only Simple Solutions, no un nivel actualmente acreditado: la revisión interna no equivale a un appraisal CMMI ni permite afirmar una calificación organizacional obtenida. ISO/IEC 25010, por su parte, organiza la evaluación del producto y no certifica el software. Las certificaciones exigidas a la empresa o al equipo se acreditarán por los mecanismos de las Bases y no sustituyen la evidencia de calidad de una versión de Ancoa.

La gestión aplicará cinco principios de trabajo seleccionados para este plan, enfoque en el CLIENTE, liderazgo, compromiso del equipo, mejora continua y decisiones basadas en evidencia. Esta selección organiza prácticas del plan y se complementa con el enfoque a procesos y la coordinación de relaciones con el CLIENTE y los proveedores. En Ancoa, el enfoque en el CLIENTE se concreta en acordar criterios y revisar evidencia con la Contraparte; el liderazgo, en que el Jefe de Proyecto y el Líder de Calidad asignen responsables y resuelvan impedimentos; el compromiso del equipo, en revisiones por pares y registro de defectos por quienes construyen y operan cada servicio; la mejora, en verificar las acciones derivadas de hallazgos al cierre de cada intervalo; y las decisiones basadas en evidencia, en conservar medición, fuente, condiciones y resultado antes de cambiar una puerta o aceptar una desviación. Los dueños de Retail y del Emisor validan por separado las reglas y los resultados de sus dominios.

En cada revisión de intervalo, Calidad y el Jefe de Proyecto comprobarán una selección documentada de prácticas de gestión de requisitos, planificación, verificación, configuración y mejora. Para cada práctica registrarán el procedimiento vigente, el ejemplo de aplicación, la revisión asociada, los hallazgos y el estado de la acción correctiva, con dueño y fecha objetivo. El resultado será una lista priorizada de brechas y acciones que se revisará en el intervalo siguiente; solo la evidencia verificada permitirá declarar implantada una práctica. Una práctica sin respaldo quedará como no verificada, y la lista no se presentará como una calificación CMMI alcanzada.

\subsection{9.1.2 Objetivos por categoría contractual}

La Tabla 9.1 relaciona los ocho títulos del Artículo 24 con medidas y condiciones de salida. Los máximos de transacción particulares del Caso se aplican según su capítulo 15 (Multitiendas Ancoa S.A., 2026, Caso 09, cap. 15, p. 34); las fichas de T-13 precisarán método, instrumento, datos, responsables y evidencia por escenario.

\textbf{Tabla 9.1: Síntesis de objetivos de calidad del producto}

\begin{longtable}{p{0.16\linewidth}p{0.21\linewidth}p{0.28\linewidth}p{0.23\linewidth}}
\toprule
Categoría & Medida & Umbral o condición & Evidencia prevista \\
\midrule
\endhead
Funcionalidad & Casos de aceptación de etapa & 100 \% aprobados y firmados & UAT y acta \\
Desempeño & Tiempo por operación & Máximos Caso y P95 común a 1,5$\times$ peak & Informe de carga \\
Compatibilidad & Flujos de integración del caso & 100 \% sin error & Contratos y pruebas \\
Usabilidad & Accesibilidad y tarea por perfil & WCAG 2.2 AA; $\leq$3 interacciones; máximos Caso & Informe manual y automático \\
Fiabilidad & Disponibilidad; RTO/RPO críticos & $\geq$99,9 \% mensual; $\leq$4 h/$\leq$15 min & Telemetría y conmutación \\
Seguridad & Hallazgos críticos/altos abiertos & 0 antes de producción & Pentest y retest independiente \\
Mantenibilidad & Cobertura; deuda; duplicación & $\geq$70 \%; $\leq$5 \%; $\leq$3 \% & Pipeline por servicio \\
Portabilidad & Escala sin rediseño; reversibilidad & $\geq$3$\times$ volumen en 3 años; ensayo conforme & Prueba y expediente de salida \\
\bottomrule
\end{longtable}

\textbf{Fuente:} Elaboración propia; máximos particulares a partir de Caso 09, cap. 15, p. 34.

La tabla concentra condiciones de producto, no acredita resultados obtenidos. Para desempeño se conservarán los máximos del Caso por transacción, sin reemplazarlos por P95. La consulta de disponibilidad digital será $\leq$400 ms, confirmación de pedido $\leq$3 s, venta completa en caja con pago externo $\leq$25 s, evaluación crediticia $\leq$8 s, propagación de precio a cajas y canal digital $\leq$5 min, consulta en sala $\leq$2 s y devolución $\leq$60 s (Multitiendas Ancoa S.A., 2026, Caso 09, cap. 15, núm. RT-09.01, p. 34). Los P95 comunes siguen siendo portal $\leq$2 s, navegación $\leq$1 s, consulta API $\leq$500 ms, escritura API $\leq$800 ms, búsqueda compuesta $\leq$3 s e informe estándar $\leq$30 s; también se comprobarán lotes $\geq$10.000 registros/min, archivo de 100 MB $\leq$60 s y arranque en frío $\leq$60 s. Cada escenario conservará reloj de inicio y fin, resultado y perfil de carga; se probarán ambos peaks del Caso por separado y combinaciones justificadas, sin convertir el total de pedidos de varios días en una tasa de llegada.

El tratamiento de usabilidad incluye prototipos y pruebas con participantes de tienda, atención, logística y Emisor antes de la construcción definitiva. Las funcionalidades principales exigirán como máximo tres interacciones desde la pantalla de inicio, conforme a RT-13.05; los tiempos de las tareas críticas conservarán los máximos del Caso. Como objetivos internos iniciales se propone una tasa de error de uso no superior a 2 % de intentos de tareas representativas y que al menos 90 % de participantes complete sin ayuda esas tareas después de hasta dos horas de instrucción y práctica guiada. Estos valores no proceden de ISO ni de las Bases; hacen verificables la tolerancia de error y la curva de aprendizaje exigidas por el Artículo 26 y permiten detectar interfaces que demanden rediseño. Un error de uso que produzca cobro indebido, exposición de datos o crédito no autorizado será bloqueante aunque el porcentaje agregado sea favorable. UX, Calidad y los dueños de proceso medirán por tarea y perfil con participantes del CLIENTE antes de aprobar los prototipos; T-13 fijará tamaño y selección de participantes, catálogo de errores, reloj y protocolo de aprendizaje. La prueba conservará intentos, errores, tiempos y participantes que completan sin ayuda; la revisión de objetivos con la Contraparte se registrará antes de construir, sin rebajar tres interacciones ni máximos contractuales. Para seguridad de funcionamiento, característica adicional de ISO/IEC 25010:2023, se analizarán pérdidas de conectividad, pagos repetidos, errores de precio y decisiones crediticias sin autorización, y se comprobará detención o degradación segura. El mapeo de usabilidad a interacción y portabilidad a flexibilidad no supone equivalencia íntegra entre ediciones.

La calidad de datos se organizará con ISO/IEC 25012:2008, comprobando completitud, exactitud y consistencia de los datos capturados y conjuntos migrados mediante validaciones, recuentos, sumas de control y reglas de negocio. Datos implementará validación en el punto de captura y publicará un tablero accesible al CLIENTE con completitud como campos obligatorios informados sobre campos obligatorios esperados, exactitud como registros contrastados correctamente contra su fuente autorizada sobre registros contrastados, y consistencia como registros que satisfacen reglas cruzadas sobre registros evaluados. Cada indicador conservará población, período y dominio; las reglas rechazarán o enviarán a remediación los datos inválidos sin incorporarlos silenciosamente al registro oficial. QA verificará rutas positivas y negativas de captura, actualización del tablero y trazabilidad de la corrección, conforme a RT-05.04. La conciliación de saldos y transacciones será exhaustiva y toda diferencia tendrá explicación aprobada. Los resultados de negocio del Anexo D se medirán por canal, instalación y dominio; en exactitud de inventario, cancelaciones o adopción, el expediente distinguirá efecto técnico de actividades compartidas con Ancoa. Las metas propuestas del Anexo D y la tolerancia de inventario se acordarán por los responsables de negocio y la Contraparte antes de la campaña correspondiente; la fecha del hito no sustituye ese acuerdo. Se mantendrá la retención de antecedentes y consentimiento del crédito durante todo su plazo y seis años más, junto con las demás retenciones del Caso (Multitiendas Ancoa S.A., 2026, Caso 09, cap. 15, núm. RT-05.10, p. 33).

\subsubsection{Criterios de medición y responsabilidades}

Cada criterio de calidad detallado en T-13 identificará, como mínimo, el requisito o decisión que lo origina; el atributo; la métrica, unidad y estadístico; la población y denominador cuando corresponda; el valor exigido y la condición de carga, ambiente, perfil y período; el método e instrumento con su versión; la frecuencia o hito; el responsable de ejecución y el revisor; la evidencia que se conservará; y la regla para tratar un resultado fuera de umbral. El informe distinguirá el dato medido, la regla de cálculo y la conclusión, y conservará la versión del producto y la configuración ensayada. Una métrica sin población o denominador definido se marcará como no comparable y no se utilizará para declarar cumplimiento.

La consulta de disponibilidad digital tendrá como atributo el tiempo de respuesta, medido en milisegundos desde la solicitud completa del cliente hasta la recepción de la respuesta completa; el máximo será 400 ms por transacción, según RT-09.01 del Caso (Multitiendas Ancoa S.A., 2026, Caso 09, cap. 15, p. 34). La prueba usará el perfil de carga declarado para el canal digital y el ambiente de Preproducción equivalente a Producción; el resultado conservará cada observación y distribución, sin convertir el máximo en un percentil. La telemetría de extremo a extremo y el informe de carga identificarán versión, configuración, perfil, ventana y transacciones medidas.

La evaluación crediticia en punto de venta tendrá como atributo el tiempo de respuesta de la evaluación, medido en segundos desde el envío de los antecedentes requeridos hasta la respuesta completa del servicio; el máximo será 8 s por operación, según RT-09.01 del Caso (Multitiendas Ancoa S.A., 2026, Caso 09, cap. 15, p. 34). Se ensayarán perfiles y datos de prueba autorizados del Emisor en Preproducción, conservando las observaciones individuales, el estado final y la trazabilidad del caso. El registro del ensayo identificará versión, ambiente, carga, denominador de solicitudes y las respuestas fallidas o incompletas, que no se excluirán del cálculo.

La propagación de cambios de precio tendrá como atributo el tiempo transcurrido desde la aceptación del cambio autorizado hasta su disponibilidad en las 380 líneas de caja y el canal digital; el máximo será 5 min para cada destino, según RT-09.01 del Caso (Multitiendas Ancoa S.A., 2026, Caso 09, cap. 15, p. 34). Se medirán por separado todos los destinos, con denominador de 380 líneas de caja y el canal digital, bajo un perfil de eventos de cambio aprobado para el ensayo. El informe conservará marcas de tiempo de origen y recepción, cantidad y denominador de destinos, excepciones y evidencia de conciliación; un promedio no ocultará un destino fuera del máximo.

La disponibilidad mensual de cada servicio crítico tendrá como métrica la proporción de tiempo en que la transacción de negocio de extremo a extremo resulta disponible durante el período mensual; el umbral mínimo será 99,9 \%, conforme al Artículo 20 de las Bases Administrativas, RT-10.01 y RT-14.03 de las Bases Transversales. El cálculo operacional se basará en la experiencia real de la persona usuaria y la transacción completa. La disponibilidad será $100\,(M-I)/M$, donde $M$ son los minutos computables del mes e $I$ los minutos de indisponibilidad; el informe identificará regla de cómputo, servicio, intervalos y exclusiones autorizadas. Las pruebas sintéticas solo complementarán el monitoreo y no sustituirán el indicador contractual, conforme a RT-14.03. La telemetría correlacionada y el informe mensual conservarán intervalos, interrupciones y fórmula; disponibilidad de infraestructura no sustituye la medición de extremo a extremo.

La recuperación de cada servicio crítico tendrá como atributos el tiempo de recuperación (RTO) y el punto de recuperación (RPO), medidos en horas y minutos desde la interrupción declarada hasta la restitución del servicio y desde la última información consistente recuperada hasta el punto de interrupción, respectivamente. Los máximos serán 4 h para RTO y 15 min para RPO, conforme a RT-07.04 de las Bases Transversales. La conmutación real documentada se ejecutará en la campaña de recuperación ante desastres; el expediente registrará marcas de tiempo, retraso de replicación, datos recuperados y conciliación de las operaciones posteriores, además del resultado y las brechas. El informe semestral conservará los valores efectivamente alcanzados y las acciones correctivas, sin declarar conformidad antes de la prueba.

La disponibilidad mensual de servicios críticos se calculará desde la telemetría de transacción completa, agregada por servicio y mes según la regla contractual; Operaciones/SRE custodiará el registro y Calidad revisará el cálculo en cada cierre mensual. La propagación de precio se medirá con marcas de tiempo de aceptación y disponibilidad por cada una de las 380 cajas y el canal digital; Integración generará el informe en cada campaña y ante cambios representativos, y Calidad cotejará destinos, denominador y excepciones. La conciliación de cartera del Emisor medirá el número de saldos y operaciones conciliados y las diferencias no explicadas sobre la población completa de cada ensayo; Datos conservará sumas de control y resultados, el dueño de negocio del Emisor aprobará la explicación por diferencia y Calidad verificará el expediente en cada ensayo completo previo al corte. Las tres métricas se reportan separadamente por servicio o dominio y se conservan con versión, período y evidencia; ninguna agrega registros de Retail y Emisor.

\subsection{9.1.3 Gobierno, responsabilidades y evidencia}

El Líder de Calidad y Pruebas mantendrá el plan, la matriz de criterios, las métricas y el expediente de cada versión, y asignará a cada criterio un responsable de ejecución y un revisor distinto del autor cuando el riesgo lo requiera. El dueño funcional o técnico propone el criterio a partir de la línea base aprobada; Calidad comprueba que indique método, ambiente, datos, resultado esperado y evidencia, y lo incorpora al T-13 antes de la campaña. Desarrollo ejecuta revisiones por pares y corrige defectos; Arquitectura resuelve criterios estructurales junto con los responsables técnicos; Seguridad verifica controles y hallazgos; Datos valida la integridad técnica de la información; los dueños de negocio de Retail y Emisor confirman las reglas y resultados de sus dominios. Calidad revisa que la evidencia respalde ambas validaciones. Ante una medición no conforme, Calidad registra el criterio, versión, evidencia, severidad y responsable; el dueño corrige y Calidad dispone la repetición del caso y la regresión afectada. La promoción permanece bloqueada cuando no se satisface una puerta aplicable o hay hallazgos críticos o altos abiertos. Los desacuerdos técnicos se elevan al Comité de Arquitectura y los impedimentos de entrega al Comité de Proyecto, con recomendación del Jefe de Proyecto; la Contraparte Técnica coordina pruebas, valida entregables y emite las conformidades de hito. Así se separan el aseguramiento de calidad (QA), que revisa la aplicación del proceso, el control de calidad (QC), que comprueba el producto, y la aceptación contractual, que corresponde al CLIENTE.

Only Simple Solutions asigna al Líder de Calidad y a dos Ingenieros QA la planificación de campañas, la ejecución de pruebas y la trazabilidad entre requisitos, resultados y defectos. El Líder define el método, revisa la suficiencia de la evidencia y decide la recomendación técnica de apertura de cada puerta; los Ingenieros QA preparan datos y casos, ejecutan pruebas, documentan hallazgos y repiten los casos corregidos. El Jefe de Proyecto reserva tiempo y recursos para corregir y volver a probar, coordina a Desarrollo, Arquitectura, Seguridad, Datos, Operación e Implantación y eleva impedimentos. Cada especialista responde por la corrección de su ámbito y entrega registros revisables a QA. La Junta de no conformidades, presidida por el Líder de Calidad, revisa desviaciones semanalmente durante la construcción y mensualmente durante la operación; las mejoras de proceso que trasciendan el contrato pasan al Comité de Calidad y Mejora Continua trimestral. Este reparto desarrolla el equipo y el gobierno interno definidos en el Subdocumento 1, secciones 1.2--1.3.

Cada medición o prueba conservará requisito relacionado, versión, ambiente, datos y condiciones de ejecución, método, resultado esperado y obtenido, evidencia, defectos y resolución. La matriz interna de SD09 organiza 115 criterios asociados a 28 resultados del Anexo D y obligaciones transversales identificadas. Es una base de trabajo parcial; no representa cobertura total del catálogo RF/RNF/RT ni acredita pruebas ejecutadas. Calidad ampliará la cadena al 100 \% de la línea base aprobada antes del cierre de cada etapa, utilizando el catálogo oficial v3.0 y conciliando cualquier cambio de versión mediante control de cambios, sin asumir la vigencia de un catálogo alternativo. T-13 y T-17 contendrán el detalle verificable y el protocolo de aceptación, respectivamente.

\begin{samepage}
La Figura 9.1 presenta el circuito de gobierno de calidad de datos. Los dueños de Retail y Emisor de Ancoa deciden las reglas y la conformidad de su dominio; Datos las ejecuta y QA verifica la evidencia. Los intercambios autorizados conservan los controles de X-01, sin reunir los registros de ambos negocios (Multitiendas Ancoa S.A., 2026, Caso 09, cap. 9, p. 25).

\begin{center}
\begin{minipage}{\linewidth}
\centering
\ossFigCaption{Gobernanza y ciclo de calidad de datos por dominio}\label{fig:9-1}
\includegraphics[width=\linewidth]{figuras/diag-09-03_gobernanza-calidad-datos.pdf}
\par\raggedright\textbf{Fuente:} Elaboración propia a partir de Caso 09, cap. 9, p. 25.
\end{minipage}
\end{center}
\end{samepage}

El circuito parte del significado y la fuente autorizada que acuerda cada dueño de datos. Datos implementa las reglas en captura y migración y publica indicadores con población, período y dominio. Ante un hallazgo, el dueño confirma su tratamiento y Datos corrige la fuente o transformación; QA repite la verificación y conserva el resultado antes de cerrar. Por ejemplo, una diferencia de existencias en Retail debe explicar su movimiento de origen, mientras una diferencia de cartera exige conciliación de saldos y operaciones del Emisor. No se compensan diferencias entre negocios ni se autoriza un cruce porque los indicadores agregados sean favorables. Si cambia un permiso o una regla de intercambio, el Jefe de Proyecto coordina la evaluación del dueño y Seguridad y aplica el control de cambios de SD06; X-01 exige finalidad, fundamento, campos mínimos, autorización y registro. El responsable de negocio confirma la conformidad de su dominio y la Contraparte Técnica recibe el expediente para la aceptación del hito correspondiente. Los incumplimientos contractuales y los cruces no autorizados mantienen cerrado el paso, aun cuando exista una solicitud de excepción.

La planificación establece criterios, métricas, listas de revisión y evidencia requerida; la gestión incorpora prevención, revisión y mejora durante el trabajo; y el control registra mediciones e inspecciones del producto y de los entregables para determinar si satisfacen los criterios aplicables. QA y QC se coordinan según el propósito de cada actividad. Una revisión puede examinar la aplicación del proceso o una característica del producto, siempre que indique objeto, responsable y evidencia. La revisión de proceso no sustituye las pruebas del producto, y una prueba favorable no demuestra por sí sola que el proceso de control de cambios se haya seguido. El diseño aprobado de una campaña, su ejecución conforme y la aceptación contractual se registran como estados separados; solo la Contraparte formaliza esta última mediante el mecanismo aplicable.

La versión del plan reunirá los estándares aplicables y su alcance, los objetivos contractuales e internos diferenciados, los roles de ejecución/revisión/aceptación, las actividades con su ubicación EDT, las métricas instrumentadas, los criterios de recepción y las herramientas de registro. Calidad comprobará estos siete elementos antes de iniciar una campaña y al modificar la línea base. El repositorio de requisitos y configuración custodiará las versiones; el pipeline guardará las decisiones automáticas; el gestor de defectos conservará el ciclo de corrección; la plataforma de observabilidad aportará series; y el expediente enlazará los registros originales para consulta del CLIENTE. Una captura de pantalla aislada no sustituye el resultado íntegro exportable ni su configuración.

El enfoque a procesos se concreta en mantener entradas, responsables y salidas de cada puerta; la gestión de relaciones se concreta en coordinar contratos de integración, ventanas y evidencias con proveedores de las plataformas que permanecen y con los responsables del CLIENTE. Calidad revisará esas dependencias junto con el Jefe de Proyecto en cada intervalo. La falta de respuesta de un tercero se registra como impedimento con impacto y acción; no se declara una interfaz probada si solo se ha ensayado su simulador.

\subsection{9.1.4 Métricas y puertas del código}

La integración continua bloqueará el cambio si la cobertura automatizada del código de lógica de negocio es inferior a 70 \%, conforme a RT-04.11 de las Bases Transversales. SD06 establece además como objetivo de calidad alcanzar al menos 80 \% en código nuevo de lógica de negocio, medido de forma independiente; este objetivo no sustituye ni compensa el mínimo contractual. Calidad y Desarrollo fijarán en T-13 la herramienta y versión, líneas o ramas ejecutables y exclusiones justificadas antes del primer pipeline. El reporte separará cobertura global de lógica de negocio y cobertura de código nuevo por servicio; cada integración conservará ambas y bloqueará resultados inferiores al mínimo contractual. La cobertura demuestra ejecución de código, no ausencia de defectos.

Adoptamos también los límites internos definidos en SD06: deuda técnica no superior a 5 \% y duplicación de código no superior a 3 \%. La razón de deuda se calculará como esfuerzo estimado de remediación respecto del esfuerzo de desarrollo del código evaluado; la duplicación declarará población, unidad y exclusiones. En cada integración, el analizador estático calculará deuda como horas de remediación estimadas divididas por horas estimadas de desarrollo del código evaluado, multiplicadas por 100. La duplicación se medirá como líneas duplicadas sobre líneas de código fuente analizadas, multiplicadas por 100; se excluirán únicamente código generado y dependencias externas, con exclusiones versionadas y revisadas por Calidad. Los reportes por servicio impedirán ocultar un incumplimiento dentro del promedio global. Calidad fijará reglas, versión del analizador y calibración del estimador de esfuerzo en T-13 antes del primer pipeline; exceder 5 \% de deuda o 3 \% de duplicación bloqueará la promoción hasta corregir y repetir. Estos límites complementan el mínimo contractual de cobertura y hacen verificable la mantenibilidad antes de cada promoción.

Fijamos como política inicial una complejidad ciclomática máxima de 15 por función o método nuevo o modificado. Mediremos cada unidad con el analizador estático del lenguaje y reglas versionadas en T-13; superar el límite bloqueará la promoción hasta dividir la unidad o aprobar una excepción acotada conforme a las condiciones indicadas más abajo. Elegimos medir por función porque el promedio del servicio podría ocultar una pieza difícil de probar; el informe mostrará también el máximo y la distribución del código completo. La cifra 15 es una decisión nuestra para dar efecto al umbral bloqueante de complejidad solicitado por el Formulario T-7, no un valor fijado por ISO ni por las Bases. Arquitectura y Calidad la calibrarán contra la línea base antes del primer pipeline; modificarla después requerirá ADR, justificación de mantenibilidad y repetición de las pruebas afectadas.

Para acoplamiento se propone bloquear cualquier ciclo entre servicios desplegables y cualquier dependencia directa no autorizada entre Retail y Emisor, con umbral de cero en ambos casos. El grafo incluirá importaciones, contratos síncronos y eventos; verificará ciclos por tipo de dependencia y contra los límites de contexto de SD03. El objetivo es conservar despliegue independiente y la frontera regulada. Arquitectura revisará además dependencias entrantes y salientes por componente y fijará límites específicos donde la línea base muestre riesgo; no se usará un máximo global arbitrario para servicios con funciones diferentes. El pipeline mantendrá el reporte del grafo y la decisión de promoción. Un ciclo entre servicios o un cruce directo no autorizado Retail/Emisor bloqueará hasta corregir, sin excepción ADR. Las excepciones internas mediante ADR se limitarán a dependencias no cíclicas y a complejidad localizada, con alcance, plazo, responsable, prueba de contrato y regresión; ninguna podrá dispensar segregación, seguridad ni exigencias contractuales. Calidad y Arquitectura formalizarán herramienta, versión, alcance y exclusiones con la Contraparte antes del primer pipeline; una puerta sin esa configuración no habilitará promociones.

La Tabla 9.2 concentra las decisiones de código que debe comprobar cada integración.

\textbf{Tabla 9.2: Métricas de código y decisión de promoción}

\begin{longtable}{p{0.22\linewidth}p{0.3\linewidth}p{0.4\linewidth}}
\toprule
Indicador & Criterio de control & Tratamiento \\
\midrule
\endhead
Cobertura automatizada de lógica de negocio & $\geq$70 \% & Mínimo contractual bloqueante en integración continua. \\
Cobertura de código nuevo de lógica de negocio & $\geq$80 \% & Objetivo adicional de SD06, medido por separado. \\
Razón de deuda técnica & $\leq$5 \% & Política SD06; fórmula y población documentadas antes del bloqueo. \\
Duplicación de código & $\leq$3 \% & Política SD06; unidad, población y exclusiones declaradas. \\
Complejidad ciclomática & $\leq$15 por función nueva/modificada & Política inicial propuesta; análisis estático y ADR. \\
Acoplamiento entre componentes & 0 ciclos entre servicios; 0 cruces directos no autorizados & Política propuesta; grafo versionado y bloqueo. \\
\bottomrule
\end{longtable}

\textbf{Fuente:} Elaboración propia.

Las métricas se interpretarán junto con el riesgo funcional y de operación. En un servicio que procese transacciones del Emisor, por ejemplo, un promedio favorable no compensará una ruta crítica sin pruebas ni una dependencia que rompa la segregación. La puerta registra la condición de avance y el motivo de cualquier excepción permitida; las obligaciones contractuales de seguridad y calidad no se dispensan mediante una decisión interna.

La deuda se registrará por origen, dominio, componente, consecuencia, responsable y fecha de resolución. Un atajo deliberado requerirá decisión documentada y trabajo de remediación; la deuda involuntaria se registrará cuando se identifique; y la obsolescencia de una dependencia exigirá revisar soporte, vulnerabilidades y pruebas de sustitución. Un defecto es un incumplimiento observable que debe corregirse según su severidad; denominarlo deuda no cambia el bloqueo. En particular, el registro no puede aplazar un cruce Retail/Emisor no autorizado, una pérdida de integridad financiera ni un hallazgo de seguridad contractual. Calidad comprobará al cierre de cada intervalo que las acciones vencidas permanezcan visibles en la planificación y que el indicador automático no oculte trabajo fuera del código.

\subsection{9.1.5 Planificación del costo de calidad}

El esfuerzo de calidad se registrará en los paquetes de la EDT y en la nivelación de recursos de T-15. El costo de calidad (COQ) se seguirá con horas imputadas por actividad, rol, etapa y paquete, separando prevención, evaluación, fallas detectadas antes de liberar y fallas posteriores a la liberación. Calidad abrirá el registro al planificar cada intervalo y lo actualizará con las horas consumidas, repeticiones, correcciones y defectos escapados; cada anotación conservará el paquete EDT y el defecto, caso o revisión que la originó. Prevención comprende revisión de requisitos y diseño, capacitación y automatización; evaluación comprende campañas, auditorías y mediciones; las fallas internas comprenden diagnóstico, corrección y repetición antes de liberar; las externas comprenden atención correctiva y restitución de datos después de liberar.

Al cierre del intervalo, el Jefe de Proyecto y el Líder de Calidad compararán horas previstas y consumidas por categoría y etapa, revisarán qué defectos escaparon y qué pruebas o controles generaron repetición, y dejarán acciones con responsable y paquete de trabajo. Si el patrón muestra retrabajo repetido, la acción deberá evaluar una medida preventiva en el siguiente intervalo; no se reducirá una prueba contractual para compensar el consumo. El informe distinguirá horas observadas de costos monetarios y no inferirá valores monetarios, ahorros ni un porcentaje objetivo del presupuesto. La reserva de 15 \% para deuda en construcción de SD06 se identificará separadamente como capacidad planificada; las horas efectivamente utilizadas se clasificarán por actividad sin sumar la reserva como esfuerzo ejecutado ni contabilizar dos veces la estabilización; la rama EDT 1.9 registra evaluación y certificación, mientras prevención y correcciones se imputan a sus paquetes de origen para evitar doble conteo. Este seguimiento informa ajustes de capacidad y estimación en T-15, sin modificar por sí mismo alcance ni hitos.

El registro de COQ conservará para cada actividad horas planificadas y reales, rol, período, paquete EDT, etapa, categoría y vínculo a la revisión, prueba, defecto o retrabajo relacionado. El Líder de Calidad contrastará las horas con las hojas de dedicación y los expedientes de campaña; las horas no vinculadas se informarán como no clasificadas hasta resolver su origen. El seguimiento permitirá estimar capacidad y reconocer causas recurrentes, no demostrar ahorro monetario ni fijar un porcentaje universal de inversión. Los importes, si se requieren para el control económico, se calcularán en la documentación económica y con las bases de cálculo correspondientes, fuera de esta sección técnica. Este análisis no autoriza reducir pruebas, seguridad, segregación, conciliación, personal ni actividades contractuales para ajustarse a una razón de costo objetivo.

Para controlar la dedicación sin fijar una cuota arbitraria, el informe calculará $E_{PE}=H_{\text{prevención}}+H_{\text{evaluación}}$ y su participación como $100\,E_{PE}/H_{\text{total}}$ sobre el mismo período y población de horas. El denominador incluirá las actividades planificadas o ejecutadas de la etapa según la vista elegida; ambos resultados se mostrarán por separado. La diferencia plan/real y las horas de fallas internas/externas orientarán la nivelación en T-15. Sin un total conciliado no se publicará un porcentaje como resultado. Las cifras ilustrativas, cuotas de industria y multiplicadores históricos de costo del cambio no son parámetros del proyecto. El propósito de la medición es reconocer capacidad y corregir causas, conservando la mejor solución y los controles obligatorios.

Por ejemplo, ante un defecto de transformación de cartera, el análisis de regla antes de construir se imputará a prevención, el ensayo y conciliación a evaluación, la corrección y repetición previas al corte a falla interna y la restitución posterior a producción a falla externa. Cada tarea se contará una vez en su paquete de origen, incluso cuando varias especialidades participen. El Jefe de Proyecto revisará si la causa exige cambiar el diseño o automatizar una comprobación adicional; la decisión se basará en integridad y riesgo, sin presentar una comparación monetaria o ahorro supuesto.

\section{9.2 Estrategia de Aseguramiento de Calidad}

La estrategia previene defectos desde el diseño y cada cambio y comprueba con evidencia que las capacidades para Retail y para la filial emisora cumplen sus criterios antes de avanzar de ambiente o de etapa. El equipo de Only Simple Solutions responde por la calidad del proceso y de los entregables; QA organiza la verificación con independencia de quienes construyen, y la Contraparte Técnica revisa y acepta los hitos mediante acta. La estrategia se aplica a las dos etapas y a la operación, se alinea con los procesos, la documentación y las técnicas de la serie ISO/IEC/IEEE 29119 (2021--2022) y deja el detalle de casos, calendario y formatos en el Formulario T-13 y el protocolo de aceptación en el T-17.

\subsection{9.2.1 Puertas, revisiones por pares y análisis de cambio}

El autor adjuntará al cambio el requisito/defecto, alcance, riesgo, análisis y pruebas. Un revisor distinto del autor comprobará criterio aprobado, contratos e interfaces, separación Retail/Emisor, rutas negativas y ausencia de datos, secretos o dependencias no autorizados. En cambios de datos revisa también el dueño del dominio; Seguridad revisa controles de seguridad y Arquitectura cambios estructurales. La lista de comprobación registra  (i) criterio vinculado; (ii) impacto en componentes, interfaces, datos y permisos; (iii) rutas positivas, negativas y de recuperación; (iv) resultados de pruebas y análisis; (v) observaciones con severidad y disposición; y (vi) aprobación fechada e identidades autor/revisor. Ningún cambio se integra sin checklist completo y aprobación distinta del autor. El dueño de negocio valida sus reglas y la Contraparte acepta los hitos; una revisión técnica no suplanta estas decisiones.

El análisis estático inspecciona código, dependencias, configuración e infraestructura sin ejecutar el sistema; el dinámico observa el comportamiento con estímulos controlados. Cada resultado conservará herramienta/versión, alcance, exclusiones, regla, artefacto, salida y vínculo al cambio. Un fallo de regla bloqueante de 9.1.4 o las Bases impide integrar/promover; una alerta no bloqueante requiere disposición y dueño antes de cerrar el cambio.

Calidad auditará la aplicación de revisión por pares, trazabilidad y puertas al cierre de cada intervalo de coordinación de SD06 y antes de las certificaciones. Revisará muestras de cambios contra registros originales, comunicará desviaciones al Jefe de Proyecto y verificará su corrección en el intervalo siguiente. El equipo usará Pareto para priorizar defectos observados, histogramas para distribuciones y diagramas causa-efecto para formular hipótesis; ninguna herramienta demuestra por sí sola una causa ni establece una proporción universal de fallas. Los incidentes críticos tendrán análisis de causa raíz y reporte dentro de cinco días hábiles, conforme al Artículo 25.

En cada revisión de intervalo, Calidad estratificará defectos por severidad, dominio, servicio, versión, origen y fase de detección. Pareto ordenará frecuencias o esfuerzo de corrección usando el mismo universo; un problema crítico de segregación se atiende aunque su frecuencia sea baja. Un histograma de latencia mostrará la distribución y la cola por operación y perfil, conservando el máximo exigido. Los límites contractuales permanecen distintos de cualquier límite estadístico de control del proceso.

Ante un defecto grave o recurrente, el responsable técnico reunirá reproducción, registros, cambios y condiciones del ambiente. El equipo utilizará causa-efecto para formular hipótesis de datos, proceso, configuración, integración o personas, y comprobará cada hipótesis con evidencia. Si pedidos aparecen duplicados, deberá distinguirse entre reintento sin deduplicación, repetición de evento o doble captura; la acción correctiva atenderá la causa demostrada y añadirá prueba que reproduzca el fallo. QA verificará la corrección y observará recurrencia en el intervalo siguiente. El informe de causa raíz identifica causa, efecto, contención, corrección, acción preventiva y responsables, dentro del plazo del Artículo 25 cuando corresponda.

La prevención combina criterios de terminado verificables, revisión por pares obligatoria, análisis estático y dinámico, control de configuración y automatización del flujo de integración. Ningún cambio llega a la rama principal sin revisión aprobada. Cada integración ejecuta compilación, pruebas unitarias, análisis estático, composición de software, detección de secretos y revisión de imágenes de contenedor. Los secretos se mantienen en un gestor con rotación automática y auditoría de acceso; las credenciales no se incorporan al código, imágenes ni archivos de configuración. DevOps y Seguridad verificarán la configuración y ensayarán la rotación automática, la renovación de los consumidores y la revocación del secreto anterior, conservando eventos de auditoría; una rotación manual no satisface ese criterio. Las fallas de compilación, pruebas o políticas bloquean la promoción; el registro conserva la revisión, la versión del código y los resultados. La definición de terminado mantiene separados los umbrales comprometidos. Se exige cobertura de lógica de negocio de al menos 70 \% como condición bloqueante, objetivo de al menos 80 \% para código nuevo, razón de deuda técnica no superior a 5 \% y duplicación de código no superior a 3 \%. La cobertura de trazabilidad llega al 100 \% de la línea base al cierre de cada etapa. La medición de duplicación declara unidad, población y exclusiones; no basta un porcentaje aislado.

\begin{samepage}
La Figura 9.2 muestra las puertas que controlan la versión candidata y el retorno ante fallas. Las primeras cuatro producen evidencia del equipo; el acta corresponde a la Contraparte Técnica.

\begin{center}
\begin{minipage}{\linewidth}
\centering
\ossFigCaption{Puertas de calidad para el paso de etapa a producción}\label{fig:9-2}
\includegraphics[width=\linewidth]{figuras/diag-09-01_puertas-calidad.pdf}
\par\raggedright\textbf{Fuente:} Elaboración propia.
\end{minipage}
\end{center}
\end{samepage}

La puerta de Cambio exige historia o defecto vinculado a requisito, revisión por pares, análisis de impacto, pruebas y aprobación del cambio. CI valida compilación y controles automatizados. QA ejecuta pruebas funcionales y no funcionales en ambiente aislado; un defecto bloqueante no resuelto impide promover. Preproducción comprueba aceptación, carga, resiliencia y el procedimiento de despliegue y reversión con topología, versiones y configuración equivalentes a producción. La puerta Acta representa la aceptación de los hitos contractuales, no una firma para cada integración ordinaria. El paso de etapa requiere expediente revisado y acta de aceptación suscrita por la Contraparte Técnica. El acta permite la autorización del paso a producción, que se ejecuta de forma automatizada, progresiva y reversible. Cada fallo retorna a la etapa responsable; tras una corrección se repite el caso fallido y la regresión afectada. Una excepción debe quedar aprobada, justificada, acotada y con responsable y fecha de cierre; no puede dispensar umbrales contractuales ni defectos críticos o altos abiertos antes de producción.

\subsection{9.2.2 Niveles, tipos y diseño de pruebas}

El plan de pruebas organiza niveles, tipos, ambientes, datos y responsables conforme a los procesos, la documentación y las técnicas de ISO/IEC/IEEE 29119. T-13 registrará planificación, diseño de casos, ejecución, incidentes y cierre; elegirá particiones equivalentes, valores límite, tablas de decisión y transiciones de estado según cada regla y su riesgo. Estas técnicas se combinarán con cobertura estructural y exploración de rutas negativas; la sola referencia a la serie no acredita cumplimiento de una campaña. En Desarrollo se aplican pruebas unitarias y de componente con datos sintéticos o anonimizados. QA cubre integración, contrato entre servicios, funcionalidad, regresión, seguridad automatizada y calidad de datos con conjuntos controlados, versionados y restablecibles. Preproducción ejecuta pruebas de sistema y aceptación, carga, estrés, resiliencia, recuperación ante desastres, accesibilidad, migración y ensayo de despliegue, con volumen representativo. La aceptación de usuario incorpora perfiles operacionales del CLIENTE y escenarios predefinidos del titular de crédito con datos protegidos y criterios observables. Producción y Recuperación ante Desastres son ambientes restringidos; sus evidencias provienen de telemetría, auditoría y ejercicios autorizados, sin trasladar datos productivos a pruebas salvo autorización y controles expresos.

En unidad y componente, Desarrollo comprobará reglas de cálculo, validaciones, transiciones e idempotencia con dependencias aisladas. Cada cambio ejecutará sus casos y los afectados por impacto; el par revisará el resultado. Una operación de reintento deberá mantener el identificador y producir un solo efecto de negocio aun cuando se reciba más de una vez. El fallo unitario o la infracción de una política bloqueante impide integrar; el reporte conserva versión, conjunto de casos y resultado, además de cobertura.

En integración, el equipo técnico y QA verificarán contratos entre servicios y con las plataformas que permanecen. Se ejercitarán mensajes válidos, inválidos, duplicados, fuera de orden, respuestas tardías y destinos no disponibles. Los contratos y simuladores estarán versionados. El simulador permite detectar defectos tempranos, pero la certificación de interoperabilidad necesita evidencia contra el servicio integrado o un mecanismo de homologación autorizado. QA repetirá la regresión de interfaces al cambiar esquemas, permisos o configuración; ninguna respuesta favorable del dominio Retail autoriza una consulta al Emisor fuera de la frontera aprobada.

En sistema, QA recorrerá los flujos completos por canal y perfil, incluidos fallos intermedios y recuperación. Para Retail, una reserva, pago, confirmación y devolución deben conservar estados e inventario consistentes; para el Emisor, evaluación, autorización y evidencia precontractual se verifican en su propio dominio. La campaña entra con una versión identificada, ambientes preparados y casos aprobados; termina con resultados por criterio, defectos y regresión revisados. Un escenario bloqueado por datos o ambiente se informa como bloqueado y no suma a casos aprobados.

El diseño de casos empleará particiones y valores límite en campos y cantidades; tablas de decisión para combinaciones de autorización, consentimiento y elegibilidad; transiciones de estado para pedidos, devoluciones y evaluación; y secuencias de reintentos y recuperación para integraciones. QA y el dueño del dominio revisarán el resultado esperado antes de ejecutar. En precio se incluirán cambios simultáneos, mensajes repetidos y destinos atrasados; en inventario, existencia agotada, reserva concurrente y dato sin confianza; en crédito, información incompleta, negativa y acceso no autorizado. Estas técnicas hacen comprobable la serie 29119 sin afirmar aplicación de cláusulas o plantillas normativas que no se hayan consultado.

La regresión automatizada de reglas y contratos se ejecutará en cada integración y antes de promover el candidato. Pruebas exploratorias, accesibilidad con usuarios, aceptación funcional y revisión de evidencia conservan intervención humana. Una prueba inestable se investigará y corregirá; su desactivación requiere registro e impide declarar cobertura del criterio afectado mientras falte una comprobación equivalente. El pipeline conservará resultados íntegros, duración, errores de infraestructura y salida de controles, con retención y acceso definidos en el plan de configuración. La aprobación corresponde a la versión ensayada; un nuevo artefacto invalida la promoción anterior hasta completar el impacto y la regresión aplicables.

\subsection{9.2.3 Ambientes, datos y migración}

Los conjuntos de prueba mantienen separados Retail y Emisor, con clasificación, minimización, anonimización cuando corresponda, control de acceso y retención definida. Las pruebas de frontera incluyen intentos de lectura, exportación o consulta cruzada no autorizada que deben denegarse; las pruebas del Emisor verifican autorización, propósito y consentimiento antes de usar datos de titulares. Se incluyen rutas negativas y de recuperación, además del flujo normal, para evitar que el éxito de una transacción o una conciliación agregada oculte una operación indebida. Cuando se aplique el muestreo de aceptación de 9.2.7, T-13 registrará población, unidad, criterio, selección aleatoria, semilla, tamaño y número de aceptación establecidos en este plan, con sus riesgos y supuestos. QA y Datos comprobarán la criticidad y la delimitación del lote antes del ensayo; las comprobaciones exhaustivas quedarán identificadas fuera de la muestra. El muestreo dirigido complementa la búsqueda de anomalías, pero no permite inferencia representativa ni sustituye conciliaciones exhaustivas de saldos o controles de segregación.

Antes de los ensayos, el Líder de Datos coordinará el perfilado de cada fuente y el saneamiento de los datos de origen. El informe identificará el defecto, registros afectados, regla incumplida y decisión de tratamiento por hallazgo; los dueños de Retail y Emisor confirmarán las correcciones o exclusiones autorizadas de su dominio. QA verificará la trazabilidad entre el informe, los cambios aplicados y el conjunto preparado para ensayar. Una decisión sobre un defecto de origen no podrá dispensar un requisito de calidad contractual. Antes de la migración definitiva, Datos ejecutará en Preproducción al menos dos ensayos completos con el alcance, reglas de transformación y secuencia de corte previstos para cada dominio; el plan incluirá la cartera viva del Emisor, sus saldos y operaciones en curso (Multitiendas Ancoa S.A., 2026, Caso 09, cap. 15, núm. RT-05.15, p. 34). En cada ensayo se medirá el tiempo total y, al terminar la carga, se contrastarán recuentos, sumas de control y una selección dirigida de registros contra las fuentes autorizadas. El Líder de Datos documentará cada diferencia con su origen, análisis, resolución y evidencia; los responsables de Retail y Emisor validarán las reglas de negocio y la conciliación de sus dominios. Toda diferencia deberá quedar explicada y confirmada por el dueño de negocio del dominio antes de continuar; Calidad verificará que la evidencia sustente esa confirmación, sin compensar diferencias entre dominios. Si una discrepancia queda sin explicar, el ensayo se declara no conforme, se detiene la decisión de corte, se corrigen datos o transformaciones y se repite el ensayo completo afectado antes de autorizar la migración definitiva. El equipo también ejecutará y registrará el procedimiento de retorno previsto, incluida la conciliación de las operaciones posteriores al corte; la Contraparte Técnica revisará el expediente y emitirá la conformidad que corresponda al hito. T-13 y T-18 detallarán los casos, evidencias y protocolo de corte y reversión.

La ejecución funcional recorre pruebas unitarias, de componente, integración, sistema y aceptación; la no funcional incluye desempeño, seguridad, compatibilidad, usabilidad y accesibilidad, fiabilidad, portabilidad y mantenibilidad. La carga se prueba en Preproducción a 1,5 veces el peak declarado y el informe presenta curvas de respuesta y carga, saturación, consumo de recursos y recuperación. Separadamente, la prueba de estrés incrementará la carga hasta identificar el punto de quiebre, con criterios de detención y recuperación segura; ese ensayo no reemplaza la aprobación de umbrales bajo 1,5 veces el peak. Los umbrales de respuesta comprometidos deben cumplirse bajo esa carga de ensayo; el informe conserva cada medición y no usa promedios o percentiles para ocultar máximos que las Bases definan como tales. La resiliencia se ensaya antes de cada paso a producción y al menos semestralmente durante la operación. La recuperación ante desastres incluye conmutación y retorno antes de cada producción y semestralmente durante la operación, con RTO y RPO medidos en conmutación real. Un tercero independiente del adjudicatario ejecuta pruebas de intrusión antes de cada paso a producción y anualmente durante la operación, y entrega íntegramente el informe y el plan de remediación. Las interfaces se revisan con herramientas automatizadas y pruebas manuales frente a WCAG 2.2 AA.

\subsection{9.2.4 Verificación y validación}



La verificación comprueba que el artefacto satisface una especificación aprobada; la validación confirma con usuarios y dueño del dominio que resuelve el uso previsto. Retail valida sus flujos comerciales; el Emisor valida por separado sus procesos, autorizaciones y resultados financieros. QA organiza y revisa evidencia; el autor no aprueba su cambio; dueños de negocio validan reglas; la Contraparte acepta hitos conforme a las Bases. Ninguna de estas decisiones se sustituye entre sí.

Para una regla de precio, la verificación cotejará la transformación y propagación con la versión autorizada, destinos y límite de tiempo; la validación hará que los responsables comerciales y usuarios de caja recorran exhibición, cobro y atención de una discrepancia. Para evaluación crediticia, la verificación comprobará autorización, respuesta dentro del máximo y conservación de la evidencia; la validación del Emisor observará que el flujo permita entregar la información y resolver sin omisiones. Un resultado rápido que elimine información obligatoria es no conforme. La decisión de negocio y la prueba técnica quedarán vinculadas, conservando evidencias y permisos propios de cada dominio.

La independencia se asigna por riesgo. Quien desarrolla no aprueba su propio cambio; QA verifica la evidencia aunque el equipo de desarrollo mantenga pruebas unitarias; la Contraparte Técnica valida entregables y suscribe actas; y el pentest corresponde a un tercero independiente. Los responsables de Retail y de la filial emisora participan en los criterios de aceptación de sus dominios. Las campañas de alto impacto, en especial consentimiento, saldos, separación de datos, continuidad y seguridad, requieren revisión de QA y del responsable de negocio antes de presentarse a la Contraparte.

\subsection{9.2.5 Trazabilidad y tratamiento de defectos}

La trazabilidad enlaza requisito o resultado, diseño, componente y configuración, cambio de código, caso de prueba, ejecución, evidencia, defecto y despliegue. Por ejemplo, para Retail, la consulta de disponibilidad digital con máximo de 400 ms se vincula con el servicio de inventario y sus reglas de conciliación, el cambio que las implementa, el escenario de consulta con existencias disponibles y agotadas, y el reporte de carga a 1,5 veces el peak antes de certificar la etapa y producir el acta correspondiente. La medición recorre la experiencia completa de consulta y conserva cada resultado, no solo el P95. Para Emisor, un criterio de consulta autorizada de cartera se enlaza con el modelo de autorización y consentimiento, las reglas de segregación, las pruebas positivas y negativas por perfil, el registro de acceso y el informe de seguridad. Ambos recorridos comprueban accesos denegados y cualquier intercambio permitido por la frontera de integración con finalidad, autorización y evidencia. Una conciliación global favorable no compensa un acceso cruzado indebido.

Cada ejecución futura generará un expediente con identificador del criterio, requisito fuente, versiones de diseño y software, ambiente y configuración, datos y perfil de carga, precondiciones y pasos, resultado esperado y observado, método e instrumento, duración y denominador cuando aplique, evidencia íntegra, defectos y responsable. El estado se registra como pendiente, conforme, no conforme o bloqueado. La repetición fallida mantiene cerrada la puerta; los defectos críticos y altos no se aceptan como excepciones de producción. La matriz vincula cada hallazgo con su resolución y evidencia nueva.

La matriz mantendrá navegables los enlaces en ambos sentidos. Se registrará un estado de enlace incompleto cuando falte cualquiera de esos elementos. Esta ausencia impide cerrar la campaña afectada; un cambio de línea base identifica enlaces y regresiones que deben reevaluarse. Las vistas de Retail, Emisor y controles compartidos autorizados permanecen separadas.

El registro de defecto contendrá identificador, dominio, requisito, versión y ambiente, pasos de reproducción, resultado esperado y observado, evidencia, impacto, severidad, responsable, estado, disposición, cambio correctivo y pruebas de confirmación y regresión. El ciclo pasa por detección, clasificación y asignación al dueño, corrección o disposición por autoridad habilitada, verificación de QA y cierre. Si la repetición falla, el defecto vuelve al responsable; QA define la regresión según el impacto y conserva la causa y el cambio correctivo. La severidad considera operación, datos, seguridad, cumplimiento y usuarios. QA revisa defectos abiertos antes de cada certificación; los bloqueantes impiden certificar y toda excepción interna conserva su responsable/plazo, sin dispensar exigencias contractuales.

\subsection{9.2.6 Criterios de aceptación de entregables}

Los tres entregables siguientes se inspeccionan antes de presentarlos al CLIENTE. Se aplican criterios documentales al plan, resultados de ejecución al paquete de migración y campañas completas a la versión candidata. Una inspección del plan no acredita la conformidad de los otros dos entregables.

\subsubsection{Aceptación del plan de pruebas y catálogo T-13}

QA verificará que todos los requisitos de la línea base aprobada para la etapa tengan método, resultado esperado, ambiente y evidencia previstos, y que cada caso identifique nivel, tipo, datos, automatización y responsable. Los requisitos y casos se revisarán por Retail y Emisor en vistas separadas. El Líder QA aprobará la integridad del plan tras resolver observaciones de revisión por pares; la Contraparte Técnica aceptará el entregable cuando corresponda a un hito. La evidencia será el plan versionado, la matriz de cobertura y el registro de observaciones resueltas. La cobertura del universo aprobado debe ser completa para cerrar el plan; una brecha no se contará como prueba aprobada.

\subsubsection{Aceptación del paquete de migración por dominio y ola}

Datos preparará el plan, perfilado, saneamiento, transformaciones, conciliación y retorno de cada ola. QA comprobará el expediente y los dueños de Retail y Emisor aprobarán por separado las reglas y resultados de su dominio. El criterio de aceptación requiere al menos dos ensayos completos en Preproducción conforme a RT-05.13, conciliación reproducible de recuentos y sumas de control, tratamiento documentado de cada diferencia y ejecución del retorno previsto. Ninguna diferencia inexplicada permite autorizar el corte. La evidencia incluye plan y bitácoras fechadas, versiones del proceso y los datos, conciliación, registro de defectos, decisión del dueño del dominio y resultado del retorno; la Contraparte Técnica acepta el hito mediante acta.

\subsubsection{Aceptación de versión candidata de etapa}

QA verificará regresión y campañas funcionales y no funcionales aplicables; Seguridad presentará la evaluación independiente prevista; Operaciones demostrará despliegue y reversión; cada dueño de negocio confirmará los resultados de su dominio. La versión candidata cumple cuando satisface los umbrales contractuales aplicables, no mantiene defectos que bloqueen la liberación conforme a las puertas definidas y presenta el expediente trazable por versión y ambiente. La Contraparte Técnica acepta la etapa mediante el acta exigida antes de producción. La evidencia comprende resultados de campaña, matriz y registro de defectos, informe independiente aplicable, despliegue/retorno y acta. Ninguna aprobación interna habilita por sí sola producción.

Antes de cada hito, QA verifica integridad y trazabilidad del expediente; el responsable de servicio declara resultados y riesgos; y el jefe de proyecto presenta el entregable a la Contraparte Técnica. Conforme al artículo 18 de las Bases Administrativas, cada presentación reúne el artefacto, evidencia objetiva, trazabilidad y registro de observaciones resueltas. La Contraparte dispone de diez días hábiles para revisar; el adjudicatario cuenta con diez días hábiles para subsanar. La reiteración de observaciones de la misma naturaleza se trata como atraso imputable. La aceptación se formaliza por acta y no libera de defectos posteriores ni convalida incumplimientos detectados después. El detalle por hito se desarrolla en T-17; el diseño de campañas y casos se registra en T-13.

\subsection{9.2.7 Muestreo y límites de la aceptación estadística}

Para atributos de registros migrados no financieros que admitan muestreo se propone un plan interno de tamaño fijo $n=299$ y número de aceptación $c=0$. La probabilidad de observar cero defectos si la proporción real alcanza 1 \% es $0{,}99^{299}\approx0{,}0495$. El cálculo representa la probabilidad de no detectar defectos bajo selección aleatoria e independencia; en extracción sin reemplazo de un lote finito, la expresión binomial se usa como cota conservadora. No significa que el lote tenga 95 \% de calidad ni que todos sus registros sean correctos. Esta regla no es umbral contractual ni estimación de calidad de Ancoa; es una política inicial del proponente, que se documentará y acordará antes de ejecutar la campaña sin rebajar controles contractuales. Si el lote contiene 299 registros o menos, se inspecciona el censo. En lotes mayores se extraen 299 registros aleatorios mediante semilla reproducible. La aceptación estadística requiere cero defectos; cualquier defecto observado rechaza el lote, inicia análisis de causa y exige censo o repetición integral de la transformación afectada. Así se limita el riesgo de aceptación del consumidor ante una tasa de defectos de 1 \% o superior; el riesgo de rechazo del productor se acepta conscientemente por  $c=0$  y se informará en T-13.

La regla se limita a atributos de registros ordinarios no cubiertos por controles exhaustivos. Los saldos y operaciones vivas del Emisor y los totales de control críticos se concilian exhaustivamente. Autorización, consentimiento y segregación requieren comprobar todas las reglas y situaciones de riesgo de la línea base; su verificación no se sustituye por esta muestra de registros. La tolerancia para diferencias financieras o de frontera es cero; cualquier discrepancia inexplicada detiene aceptación y corte. El muestreo dirigido busca anomalías raras/extremas, pero no estima prevalencia. Retail y Emisor conservan poblaciones, propietarios, reglas y decisiones separadas. Toda corrección sistémica requiere repetir la transformación completa afectada.

Cada lote se delimitará por dominio, fuente, ola, transformación y criterio de atributo, con fecha de extracción y listado íntegro de unidades. Datos registrará la semilla y QA verificará selección uniforme, ausencia de duplicados y trazabilidad de las unidades elegidas. Los estratos de alto riesgo se examinarán separadamente; una selección dirigida adicional no se contará dentro de la muestra aleatoria ni cambiará su interpretación. Si se rechaza el lote, se aislará su carga, se corregirá la causa y se repetirá la transformación y la inspección antes de volver a decidir.

Con $c=0$, el riesgo de rechazo de un lote con proporción real $p$ se aproxima por $1-(1-p)^{299}$ bajo el modelo binomial. Para un escenario hipotético de $p=0{,}001$, el riesgo es aproximadamente 25,9 \%; este cálculo expone la severidad del plan y no estima defectos de Ancoa. QA informará conjuntamente tamaño, límite, supuestos y riesgos antes de la decisión de la Contraparte. Un resultado sin defectos en la muestra permite la decisión estadística acotada al criterio y lote; no valida otros atributos, conciliación financiera, derechos de acceso ni calidad global de una etapa.

\section{9.3 Alineación con Plan de Trabajo}

La alineación asigna productos de calidad a las ramas de trabajo y protege su ejecución durante las certificaciones, el solapamiento y la operación. Cada decisión de avance conserva la versión, el expediente y la autoridad que corresponde al hito contractual.

\subsection{9.3.1 Productos de calidad, EDT e hitos}

La calidad se programa como trabajo verificable dentro de la EDT y acompaña construcción, migración, marchas blancas, aceptación y operación. La rama 1.9 reúne el plan de pruebas, revisiones, baterías funcionales y no funcionales e informes de certificación de ambas etapas. La rama 1.8 aporta el modelado de amenazas, matriz de controles, segregación de funciones e informes de pruebas de intrusión. La rama 1.7 produce planes, ensayos y actas de migración, incluida la cartera viva del Emisor. La rama 1.12 organiza los planes e informes de marcha blanca y las actas de aceptación por etapa. Los productos y dependencias de estas ramas se coordinan con el cronograma y la EDT del Capítulo 7, sin sustituir el Gantt ni alterar sus hitos contractuales.

\begin{samepage}
La Figura 9.3 ubica las certificaciones y marchas blancas de ambas etapas. La Etapa 2 comienza durante la marcha blanca de la primera y la operación contractual continúa hasta M56.

\begin{center}
\begin{minipage}{\linewidth}
\centering
\ossFigCaption{Calidad sobre los hitos del proyecto}\label{fig:9-3}
\includegraphics[width=\linewidth]{figuras/diag-09-02_hitos-calidad.pdf}
\par\raggedright\textbf{Fuente:} Elaboración propia.
\end{minipage}
\end{center}
\end{samepage}

La figura muestra dos ciclos de calidad que comparten plataformas y controles, pero conservan criterios de aceptación por etapa. La certificación del mes 12 prepara los entregables y campañas de la Etapa 1; en los meses 13 a 15 se recopila evidencia de la marcha blanca y se corrigen hallazgos antes del acta y la producción del mes 16. La certificación del mes 18 prepara el segundo corte mientras la Etapa 1 ya está en producción; los meses 19 y 20 corresponden a la marcha blanca de Etapa 2 y a su convivencia controlada con la primera etapa. El acta habilita la producción del mes 21, cuando comienza el período de operación contractual hasta el mes 56.

En la rama 1.9, QA mantiene casos y resultados vinculados a la línea base, mide umbrales y prepara las certificaciones de los meses 12 y 18. En 1.8, Seguridad verifica que los controles y las correcciones se implementen antes de esas certificaciones, y coordina el pentest independiente previo a cada liberación. En 1.7, el equipo de datos y los responsables de negocio ejecutan los ensayos de migración y documentan conciliación, excepciones, corte y retorno antes de la aceptación que dependa de ellos. En 1.12, Implantación y QA operan las marchas blancas, conservan mediciones y observaciones y preparan las actas. La Contraparte Técnica revisa y acepta los hitos; la filial emisora valida los criterios que afectan cartera, autorización y controles financieros.

Durante la marcha blanca, la solución opera con datos y usuarios reales en paralelo con la operación vigente, sin ser aún el sistema oficial de registro. Los responsables de cada etapa comparan resultados, vigilan incidencias y aplican criterios de detención y retorno ya ensayados. El cierre exige conjuntamente cero incidentes críticos o altos atribuibles a la solución, volumen real comprometido durante al menos las últimas cuatro semanas, disponibilidad y tiempos sostenidos en ese período, conciliación sin diferencias no explicadas, personal capacitado y certificado según el plan aprobado y acta firmada por la Contraparte. Si las condiciones de salida no se satisfacen, se amplía la marcha blanca a cargo del adjudicatario, sin cargo adicional y sin desplazar las fechas contractuales de fases posteriores; se atienden las consecuencias por atraso establecidas en las Bases. La fecha objetivo de producción no sustituye la evidencia ni habilita a omitir observaciones, acta o controles de seguridad.

Desde el mes 21, la rama 1.15 asume la calidad en operación y mantiene los procedimientos establecidos en 1.9 y los controles de 1.8. Esto incluye seguimiento de defectos y regresiones, controles de seguridad, pruebas de resiliencia y recuperación con la frecuencia comprometida, y pruebas de intrusión anuales. Los responsables de servicio y Operaciones conservan indicadores y evidencia de despliegues; QA actualiza la matriz de trazabilidad cuando cambia un requisito, configuración o componente. La Contraparte verifica los compromisos operacionales en los hitos y revisiones aplicables. Las pruebas de migración y conciliación cerradas en 1.7 siguen respaldando los datos operativos, y 1.12 registra estabilización y aceptación final.

La preparación de calidad comienza con criterios, plan de pruebas, reglas de revisión y ambientes identificados; estos productos habilitan construir y medir, y no se posponen hasta la certificación. En 1.9 el responsable QA mantiene el catálogo de casos, ejecuciones e informes de etapa; la evidencia de cobertura y bloqueos acompaña cada candidato. En 1.8 Seguridad aporta la matriz de controles, hallazgos y remediación verificada; en 1.7 Datos aporta perfilado, transformaciones, ensayos y conciliación por dominio. La aceptación requiere reunir estas contribuciones, sin declarar que cada rama sea un hito económico autónomo.

La implantación y marcha blanca de 1.12 incorporan las bitácoras de despliegue y retorno, series de operación supervisada, conciliación, observaciones y evidencia de formación exigida. El responsable de cada rama identifica en T-14 el paquete que produce cada registro y sus dependencias. El Jefe de Proyecto comprueba que el expediente de certificación llegue antes de la ventana de revisión y que los defectos tengan trabajo de corrección y repetición asignado. Si falta un insumo habilitante, el calendario interno debe recuperar esa dependencia; el hito contractual conserva su fecha y su condición de aceptación.

\subsection{9.3.2 Capacidad, ambientes y versiones durante el solapamiento}

La coordinación de calidad en los meses 13 a 20 se trata como dos corrientes de trabajo con prioridades y líneas base separadas. Se atienden el cierre y operación supervisada de la Etapa 1, y construcción y certificación de la Etapa 2.

Antes de cada ventana de certificación, el Líder de Calidad y Pruebas contrasta la curva mensual y por frente del Formulario T-15 con el calendario de campañas de T-13 para ejecución funcional y no funcional, revisión independiente, triage y repetición de pruebas, preparación del expediente y atención de observaciones. Líderes de Desarrollo, Seguridad, Datos, Implantación y Operación confirman disponibilidad de las personas que producen o validan evidencia. La asignación preserva la independencia ya establecida en 9.2. Quien desarrolla no aprueba su propio cambio y las pruebas de intrusión las ejecuta un tercero independiente. Si dos campañas requieren simultáneamente un mismo recurso crítico, el Jefe de Proyecto coordina la nivelación antes de fijar la fecha interna del gate y registra impacto, decisión y recuperación en el plan integrado; no declara capacidad disponible sin respaldo de T-15.

Los ambientes de Desarrollo, QA, Preproducción, Producción y Recuperación ante Desastres se mantienen identificados por inventario de componentes, topología, configuración, permisos y versión. Desarrollo usa datos sintéticos o anonimizados; QA utiliza conjuntos controlados y versionados; Producción y DR mantienen acceso restringido y auditado. Durante el solapamiento se registra por separado el propósito y ventana de uso de cada ambiente, los cambios autorizados y la versión de datos sintéticos o enmascarados, de manera que una prueba de E2 no modifique datos, configuración ni telemetría de operación de E1.

Cada candidato queda identificado por una versión inmutable del código y del artefacto desplegable, con dependencias, imágenes, configuración externalizada y resultados de pipeline asociados. La promoción a Preproducción y Producción reutiliza el mismo artefacto aprobado, sin recompilar; solo cambia la configuración declarada del ambiente. El expediente registra hash o identificador equivalente, versiones de componentes y herramientas que afecten la medición, parámetros y resultados. Las versiones concretas de herramientas, navegadores, motores y dependencias se fijan en la línea base de configuración aprobada para la campaña, no se presuponen en este plan. El control de configuración y la separación entre despliegue y liberación permiten observar E1 en producción y validar E2 sin exponer funciones no aceptadas; los calendarios de congelamiento de SD07 siguen rigiendo las ventanas operativas.

La evidencia de carga separa los perfiles presencial y digital y conserva cada perfil del Caso, su volumen, concurrencia y tasa de llegada definida para el ensayo. La concurrencia temporal mantiene los umbrales y perfiles diferenciados. T-13 identificará campaña, ambiente, versión y recursos compartidos para interpretar cada resultado. Antes de un gate se adjunta la comparación entre perfil observado, perfil requerido y configuración usada; una desviación de ambiente o capacidad requiere registrar su impacto y repetir la prueba en condiciones válidas cuando el criterio no pueda demostrarse con la evidencia disponible.

El calendario reserva las ventanas del artículo 18 para revisar cada entregable y subsanar observaciones antes de la compuerta correspondiente. Para ello, QA entrega el expediente completo con anticipación suficiente y el jefe de proyecto incorpora revisión y corrección en la planificación de cada corte. Los informes operacionales incluirán proyección trimestral de crecimiento, alertas anticipadas y propuesta de ajuste de capacidad, con perfil y supuestos actualizados. Esta coordinación reduce el riesgo de repetir observaciones y conserva las fechas de producción de los meses 16 y 21 y el término contractual en el mes 56. Las dependencias de infraestructura, datos, seguridad, formación y disponibilidad de usuarios se controlan en el plan integrado de SD07; cualquier cambio que afecte alcance, riesgo, evidencia o fecha pasa por el control de cambios de SD06 y requiere evaluar de nuevo las pruebas afectadas.

\subsection{9.3.3 Gestión de cambios y aceptación contractual}

Todo cambio de requisito aprobado, diseño, datos, interfaz, código, dependencia, configuración o ambiente que pueda alterar un criterio abre una evaluación de impacto antes de incorporarse. El solicitante vincula la petición o defecto con requisito, servicio, etapa y paquete EDT; el responsable técnico analiza los componentes y versiones afectados; QA determina casos, datos, criterios y evidencia que deben actualizarse o repetirse; Seguridad y los dueños de Retail o Emisor revisan los impactos de su dominio. El Jefe de Proyecto coordina la decisión por el procedimiento de control de cambios del SD06 y actualiza línea base, EDT, calendario y trazabilidad según corresponda. Una modificación aprobada no se considera verificada hasta que el candidato resultante completa las pruebas afectadas y la regresión definida. Si cambia el alcance o una condición de aceptación contractual, se aplica el procedimiento y la autoridad que establecen las Bases y el contrato; este plan no amplía unilateralmente esos compromisos.

La cadena de recepción vincula requisito, medición, prueba y evidencia con el paquete EDT y el hito de E-25. El acta acredita el cumplimiento del hito para el tratamiento contractual del pago; este capítulo no reproduce porcentajes ni precios. Una aprobación técnica automatizada o un informe del proponente no sustituye el acta. La aceptación mantiene las responsabilidades por defectos posteriores previstas en el Artículo 18.4.
Para los criterios de disponibilidad digital, evaluación crediticia y migración de cartera, la cadena comienza en el límite o regla fuente; continúa en la ficha de métrica, caso y resultado por versión; y reúne el informe de carga o conciliación con la contribución de la rama 1.9 o 1.7 al expediente de etapa. El expediente se presenta al hito que E-25 define para el entregable correspondiente y termina en el acta de la Contraparte. QA verifica los enlaces en ambos sentidos. La conformidad de cartera incorpora aprobación del dueño del Emisor, mientras los resultados Retail conservan su propio dueño y evidencia. El porcentaje y monto del estado de pago pertenecen a la oferta económica.

Operaciones entregará el informe mensual de niveles de servicio dentro de los primeros cinco días hábiles del mes, con indicadores, incidentes, causas y acciones conforme al Artículo 79. La evidencia estará disponible y será exportable para el CLIENTE. El registro distinguirá ventana aprobada, exceso de ventana, período, transacción, atribución e impacto; una discrepancia de medición se tramitará con sus datos originales y justificación metodológica. La Contraparte y el Administrador del Contrato aplican el procedimiento del Artículo 80 cuando corresponda, incluida notificación fundada, descargos y resolución. Calidad aporta evidencia y seguimiento de corrección, sin crear sanciones ni reconstruir fórmulas o ponderaciones ilegibles de las tablas extraídas.

La marcha blanca y recepción se rigen por los Artículos 17 y 18 de las Bases Administrativas, y la aceptación final del mes 21 incorpora la Garantía de Correcto Funcionamiento del Artículo 37, con vigencia de veinticuatro meses. Los defectos posteriores imputables a diseño, construcción o implantación se registrarán, corregirán y volverán a verificar con trazabilidad a la versión afectada. El cierre documental no elimina esta responsabilidad. Los procedimientos de soporte, garantía y reversibilidad conservarán los registros necesarios después de la aceptación y durante el período aplicable, sin trasladar a Ancoa una recepción provisional, retención o plazo tomado de ejemplos externos.

\subsection{9.3.4 Continuidad del control durante la operación}

Desde el mes 21, Operación/SRE actualiza el expediente al promover cada versión y conserva telemetría, incidentes, cambios, defectos, regresiones y medidas correctivas vinculadas con los mismos criterios y servicios. QA verifica la integridad de la trazabilidad en las revisiones previstas; Seguridad y el tercero independiente entregan las evidencias de intrusión con la frecuencia comprometida, y Operación registra las pruebas semestrales de resiliencia y recuperación ante desastres. Un cambio operacional que afecte un control o umbral vuelve al análisis de impacto y a las pruebas correspondientes antes de su liberación. Los informes y evidencias se relacionan con los hitos mensuales de operación que fija E-25; la aceptación y el pago se rigen por ese instrumento y por el contrato durante los meses 21 a 56.

\section*{Referencias}
\addcontentsline{toc}{section}{Referencias}

Multitiendas Ancoa S.A. (2026). \textit{Caso 09. Cadena Multitienda: Bases Técnicas del Caso}. \url{https://drive.google.com/file/d/1cVj84Af4bNYeXu8R1FHi8BoYtOLfcGs1/view}

\section*{Declaración de uso de IA}
\addcontentsline{toc}{section}{Declaración de uso de IA}

Se utilizó OpenAI Codex para contrastar fuentes, organizar y redactar este capítulo, elaborar sus figuras y revisar su coherencia documental. La intervención en texto y diagramas fue alta. Se aplicó humanizer para eliminar patrones de redacción automatizada, manteniendo el registro formal, cifras, títulos e identificadores. Esta asistencia no acredita ejecución de pruebas, certificación ni aceptación contractual. La elaboración y revisión humana exigidas por el Comunicado 10 deberán registrarse antes de presentar la oferta; no se atribuye aquí una revisión humana realizada. El uso se incorporará al Formulario A-6.

\ossFinalPage
\end{document}
